% =============================================================================
\chapter{Considerações finais}\label{cap:consideracoes}
% =============================================================================

O Projeto Analítico A3 transformou uma base pública de 27,9 milhões de negócios
de BTC/USDT em evidências para as decisões de uma mesa de execução de
criptoativos, cobrindo todas as etapas previstas nas diretrizes da disciplina:
diagnóstico, ETL com auditoria de qualidade, análise exploratória, modelagem
preditiva, visualização e reflexão ética.

A primeira contribuição do trabalho está na qualidade dos dados. Uma base que
parecia limpa, sem nenhum nulo, escondia um carimbo de ingestão 100\,\% inválido
e 41,7\,\% de duplicatas aparentes que, se removidas por um procedimento
ingênuo, teriam destruído quase metade dos negócios. A reconciliação minuto a
minuto com a fonte oficial, com concordância total nos 10.080 minutos, deu
segurança para todas as análises seguintes e é a melhor ilustração do princípio
\textit{Garbage In, Garbage Out} no projeto.

As três perguntas de negócio foram respondidas com números. Na P1, a liquidez se
concentra entre 14h e 16h~UTC, com pico 4,2 vezes maior que o vale, mas acompanhada
de volatilidade mais alta, o que recomenda execução fatiada nessa janela. Na P2,
os dias de estresse multiplicam o volume por 2,7 e a volatilidade por 2,0, e o
valor negociado é dominado por poucos negócios grandes, o que exige limites de
risco dinâmicos e folga de capacidade. Na P3, o desequilíbrio do fluxo explica o
preço do mesmo minuto, mas não o do minuto seguinte.

A modelagem preditiva confirmou e refinou essas conclusões. A volatilidade do
intervalo seguinte é previsível fora da amostra, e a regressão Ridge superou
tanto os \textit{baselines} quanto os modelos de árvore. A direção do preço, por
outro lado, não oferece ganho que cubra os custos de transação, nem mesmo na
escala de 1 segundo, em que a alta acurácia apenas reflete o \textit{bid-ask
bounce}. Para a mesa, o valor prático está na gestão de risco e no ritmo de
execução, e não em apostas direcionais.

As conclusões devem ser lidas à luz das limitações discutidas no
\autoref{cap:etica}: uma única semana, atípica, de uma única corretora e de um
único par. Como trabalhos futuros, recomenda-se:

\begin{alineas}
  \item repetir as análises em mais semanas, incluindo períodos calmos e de alta,
        e em outras corretoras e pares, para validar a estabilidade dos padrões
        horários e de estresse;
  \item incorporar dados do livro de ofertas, para medir o \textit{spread}
        cotado, a profundidade e o OFI na forma original de
        \citeonline{cont2014};
  \item comparar a previsão de volatilidade com modelos econométricos clássicos
        de volatilidade realizada e avaliar seu uso em uma simulação de execução
        TWAP e VWAP com custos;
  \item evoluir o pipeline para processamento contínuo, com o mesmo contrato de
        qualidade aplicado a cada novo lote de dados.
\end{alineas}

No momento desta documentação, as pendências do projeto, registradas em
\code{TODO.md}, são a elaboração dos \textit{slides} da defesa final, a
atualização da tabela de entregas no \code{README.md} e uma revisão de
reprodutibilidade de ponta a ponta dos notebooks, do código consolidado e do
painel antes da entrega.
